\section{实战案例：机器人控制与多智能体}

\subsection{Robotics 中的 PPO/SAC 选择}

\begin{figure}[H]
\centering
\includegraphics[width=0.9\textwidth]{slide-18.jpg}
\caption{Lecture slide 18 展示的 robotics benchmark：PPO 用于稳定 on-policy 控制，SAC 擅长高维连续动作。}
\end{figure}

在现实机器人控制任务中，训练通常面临采样限制与高维动作空间。PPO 提供的 trust-region + clipping 架构便于在真实设备上限制步幅，而 SAC 的 entropy-optimal off-policy buffer 让拟合速度更快。

\begin{knowledgebox}{现场调度的建议}
在 robot arm / locomotion 任务中常用策略是：用短 horizon 的 PPO trial 快速验证 reward 结构；将最有 promise 的设定迁移到 SAC，配合 soft target 和 automatic entropy tuning。PPO 结果更好解释、SAC 更节省样本。
\end{knowledgebox}

\subsection{多智能体训练中的共享 critic}

当多个 agent 共享环境时，common critic 或 centralized value 是提高稳定性的关键。通常做法是训练一个 joint critic $Q(\boldsymbol{s}, \boldsymbol{a})$，再把 actor 分段更新：

$$\hat{A}_i = Q(\boldsymbol{s}, \boldsymbol{a}) - V(\boldsymbol{s})$$

Agent 只观察自身 action，但 critic 看到全局状态，能够更准确判别局部决策的价值。

\begin{warningbox}{共享 critic 的风险}
如果 critic 过度依赖 global information，或没有合适的 value decomposition，单个 agent 的 gradient 会被全局 noise 掩盖。解决方案包括 information bottleneck、agent-specific mask 与 periodic critic reset。
\end{warningbox}

\subsection{本章小结}

PPO 与 SAC 在智能体训练中扮演不同角色：PPO 当作安全策略，SAC 管理高频采样；多智能体场景则需共享 critic 并注意信息泄露。

